\documentclass[11pt,a4paper]{article}
\usepackage[margin=1in]{geometry}
\usepackage{amsmath,amssymb,amsthm}
\usepackage{hyperref}
\usepackage{booktabs}
\usepackage{listings}
\usepackage{xcolor}

\lstset{
  basicstyle=\ttfamily\small,
  breaklines=true,
  frame=single,
  language=C++
}

\newtheorem{theorem}{Theorem}
\newtheorem{definition}{Definition}
\newtheorem{lemma}{Lemma}
\newtheorem{remark}{Remark}

\title{Cross-Implementation Verification of Cryptographic Primitives\\
Through Bilingual Golden-Vector Oracle Systems}

\author{Sk Sarin Arfin\thanks{Independent researcher. Contact: \texttt{sarinsk629@gmail.com}}}

\date{September 2026}

\begin{abstract}
We present a methodology for verifying the correctness of cryptographic
implementations through bilingual cross-examination: a Python oracle
independently mirrors every C++ operation, golden vectors pin every
intermediate value, and an append-only ledger owns every defect as a
numbered erratum. We apply this methodology to a complete cryptographic
proof system---elliptic curve arithmetic over the Pasta cycle
($\mathbb{F}_p$ and $\mathbb{F}_q$), Poseidon hashing, Sparse Merkle
trees, threshold BLS signatures, commitment schemes, and a folding-based
recursive state transition proof---comprising 41 implementation modules
and 34 automated verification gates. The methodology caught 162 defects
across multiple defect classes, including two classes (out-of-bounds
memory access and incorrect constraint algebra) that were structurally
invisible to the implementation's own test suite. We formalize the
oracle-divergence detection guarantee, prove the golden-vector coverage
completeness for arithmetic operations, and demonstrate the methodology's
scalability: the resulting system produces a 1{,}600-byte recursive
state transition proof ($\pi_E$) verifiable without the witness in
118\,ms on ARM64 hardware.
\end{abstract}

\section{Introduction}

\subsection{The Verification Problem}

Cryptographic implementations are notoriously difficult to verify.
A single incorrect operation---a wrong Montgomery reduction constant,
a misidentified bit ordering, an omitted domain-separation
tag---can silently compromise the security of the entire system.
Traditional approaches rely on:
\begin{enumerate}
\item \textbf{Test vectors}: fixed input-output pairs from reference
implementations. These verify specific values but do not guarantee
correctness for all inputs.
\item \textbf{Formal verification}: tools such as EasyCrypt or
F* that produce machine-checked proofs. These provide the strongest
guarantees but require significant expertise and effort.
\item \textbf{Code review}: manual inspection by experts. This is
error-prone and does not scale.
\end{enumerate}

We propose a fourth approach: \emph{bilingual cross-examination}.
Two independent implementations---one in C++, one in Python---compute
the same cryptographic operations. A golden vector pins every
intermediate value. When the two implementations disagree, the
divergence is detected immediately. When they agree across thousands
of test cases, the confidence in correctness is substantially higher
than either implementation alone.

\subsection{Our Contribution}

We make the following contributions:

\begin{enumerate}
\item We formalize the \emph{golden-oracle methodology} for verifying
cryptographic implementations (Section~\ref{sec:methodology}),
including the oracle-divergence detection guarantee
(Section~\ref{sec:divergence}).
\item We present three defect case studies that demonstrate the
methodology catching defects invisible to conventional testing
(Section~\ref{sec:casestudies}): an out-of-bounds memory access
invisible to self-consistency checks, a wrong-power invariant checker
masked by degenerate test inputs, and a bit-order convention mismatch
detected by an orientation probe.
\item We demonstrate the methodology's scalability by applying it to
a complete cryptographic proof system (Section~\ref{sec:system})
comprising 41 modules spanning elliptic curve arithmetic, hash
functions, commitment schemes, and folding-based recursive proofs.
\item We prove the completeness of golden-vector coverage for
arithmetic operations (Section~\ref{sec:completeness}).
\item We release the complete implementation, the decision ledger
(234 decisions, 163 laws), and all 34 verification gates as open
source (Section~\ref{sec:artifacts}).
\end{enumerate}

\section{Preliminaries}
\label{sec:preliminaries}

\subsection{Notation}

We work over two prime-order fields:
\begin{align*}
\mathbb{F}_p &: p = 2^{254} + \texttt{0x224698fc0994a8dd}\cdot 2^{192} + \cdots \\
\mathbb{F}_q &: q = 2^{254} + \texttt{0x224698fc094cf91b}\cdot 2^{192} + \cdots
\end{align*}
These are the scalar and base fields of the Pasta curve cycle:
Pallas ($y^2 = x^3 + 5$ over $\mathbb{F}_p$, group order $q$) and
Vesta ($y^2 = x^3 + 5$ over $\mathbb{F}_q$, group order $p$).

We write $\mathsf{Mont}(x) = x \cdot R \bmod p$ for the Montgomery
representation, where $R = 2^{256} \bmod p$, and
 $\mathsf{Canon}(x) = x \cdot R^{-1} \bmod p$ for the canonical form.

\subsection{The Pasta Curve Cycle}

The Pasta cycle consists of two elliptic curves whose base and scalar
fields are swapped:
\begin{align*}
\textsf{Pallas}&: y^2 = x^3 + 5 \text{ over } \mathbb{F}_p, \quad \#\mathsf{Pallas}(\mathbb{F}_p) = q \\
\textsf{Vesta}&: y^2 = x^3 + 5 \text{ over } \mathbb{F}_q, \quad \#\mathsf{Vesta}(\mathbb{F}_q) = p
\end{align*}
This cycle natively supports recursive proof composition without
foreign-field overhead: arithmetic in one curve's scalar field is
native arithmetic in the other curve's base field.

\subsection{The Poseidon Hash}

Poseidon-3 is a sponge-based hash over $\mathbb{F}_p$ with
 $t=3$ lanes, $\alpha=5$, $R_F=8$ full rounds, and $R_P=56$ partial
rounds. The permutation operates on a 3-element state
 $(s_0, s_1, s_2)$ where $s_2$ is initialized with a domain-separated
initial vector.

\section{The Golden-Oracle Methodology}
\label{sec:methodology}

\subsection{Overview}

The methodology consists of four components operating as a verification
pipeline:

\begin{enumerate}
\item \textbf{The Python Oracle}: an independent implementation of
every cryptographic primitive in Python, using arbitrary-precision
integers (no Montgomery arithmetic, no fixed-width limitations).
The oracle serves as the mathematical reference.

\item \textbf{The C++ Implementation}: the production implementation
using fixed-width Montgomery arithmetic for performance.

\item \textbf{Golden Vectors}: emitted by the Python oracle, containing
every intermediate value for a set of test inputs. These vectors pin
the expected output at each step.

\item \textbf{The Gate}: an automated verification pipeline that:
(a) regenerates all golden vectors from the oracle,
(b) compiles all C++ modules from source,
(c) runs all conformance tests against the fresh goldens,
(d) aborts on any failure.
\end{enumerate}

\subsection{The Golden Vector}

\begin{definition}[Golden Vector]
\label{def:golden}
A golden vector for a cryptographic operation $f$ with inputs
 $x_1, \ldots, x_n$ is a tuple
 $(x_1, \ldots, x_n, y_1, \ldots, y_m)$ where $y_i = f_i(x_1, \ldots, x_n)$ for each output $f_i$, computed by the Python oracle using arbitrary-precision
integers.
\end{definition}

The golden vector is \emph{authoritative}: the C++ implementation must
produce bit-identical output. Any divergence indicates a bug in the C++
implementation (or, less commonly, in the oracle itself — in which case
the oracle is corrected and the vector regenerated).

\subsection{The Append-Only Ledger}

Every defect discovered during development is recorded as a numbered
erratum (CA-R$nnn$) in an append-only decision ledger. Each entry contains:
\begin{itemize}
\item The defect description
\item The evidence that convicted it (probe output, golden mismatch, etc.)
\item The fix
\item A law: a generalizable rule that prevents recurrence
\end{itemize}

As of this writing, the ledger contains 234 decisions (DEC-001 through
DEC-235) and 163 laws (CA-R001 through CA-R163).

\subsection{The Gate}

The verification gate is a shell script that executes the following
pipeline:

\begin{enumerate}
\item \textbf{Clean-first}: remove all previous build artifacts
(binaries, object files, build caches)
\item \textbf{Regenerate}: execute the golden-vector emitter
(Python) to produce all 41 header files
\item \textbf{Configure}: CMake configure with Ninja
\item \textbf{Build}: compile all test binaries (abort on failure)
\item \textbf{Freshness check}: verify every binary is newer than
every golden header it depends on
\item \textbf{Diagnostics}: run structural diagnostics
\item \textbf{Conformance}: run all 34 test binaries
\item \textbf{Verdict}: GATE GREEN if and only if all steps pass
\end{enumerate}

\subsection{The Cross-Implementation Principle}

\begin{definition}[Bilingual Cross-Examination]
Two implementations $I_1$ (C++) and $I_2$ (Python) are
\emph{cross-examined} if for every operation $f$ and test input
 $\vec{x}$, both implementations compute $f(\vec{x})$ and the
results are compared. A \emph{divergence} is any difference in output.
\end{definition}

The key property is that $I_1$ and $I_2$ share \emph{no code path}.
They are written in different languages, use different arithmetic
models (fixed-width Montgomery vs.\ arbitrary-precision), and have
different optimization characteristics. A bug in $I_1$ that causes
incorrect output will be detected by comparison with $I_2$ (and
vice versa), unless the same bug exists in both implementations.

\section{Oracle-Divergence Detection}
\label{sec:divergence}

\subsection{The Divergence Guarantee}

\begin{theorem}[Divergence Detection]
\label{thm:divergence}
Let $I_1$ and $I_2$ be two implementations of a cryptographic operation
 $f$ over domain $D$. Let $T \subseteq D$ be the set of test inputs for
which golden vectors are generated and cross-checked. If $I_1(x) \neq
I_2(x)$ for some $x \in D$, and $I_1$ and $I_2$ share no code path,
then for any $x \in T$ where $I_1(x) \neq I_2(x)$, the divergence is
detected with probability 1.
\end{theorem}

\begin{proof}
By construction: the golden vector pins the output of $I_2$ (the
oracle) for each $x \in T$. The C++ implementation $I_1$ computes
 $f(x)$ independently. The gate compares the two outputs. If they
differ, the gate reports FAIL and aborts. There is no code path
where a divergence on a tested input is not detected, because the
comparison is unconditional.
\end{proof}

\begin{remark}
The guarantee is trivially strong for tested inputs. The interesting
question is: how well does $T$ cover $D$? For arithmetic operations,
 $T$ can be chosen to cover the edge cases (zero, one, $p-1$, carry
boundaries) exhaustively. For permutation-based constructions
(e.g., Poseidon), $T$ can include all intermediate states, providing
coverage of the internal computation.
\end{remark}

\subsection{Golden-Vector Coverage Completeness}

\begin{theorem}[Arithmetic Coverage]
\label{thm:coverage}
For any binary field operation $\circ: \mathbb{F}_p \times \mathbb{F}_p
\to \mathbb{F}_p$ implemented as $N$ distinct code paths (e.g.,
carry/no-carry, overflow/no-overflow, sign-difference), there exists
a test set $T$ of size at most $2N$ such that $T$ exercises every
code path.
\end{theorem}

\begin{proof}
Each code path is triggered by a specific condition on the operands
(e.g., $a_i + b_i \geq 2^{64}$ triggers a carry in limb $i$).
For each condition, we construct two test inputs: one that triggers
the condition and one that does not. The union of these inputs covers
all code paths. Since there are at most $N$ conditions, the test set
has at most $2N$ elements.
\end{proof}

\begin{remark}
In practice, we generate 516 golden vectors for the field arithmetic
(covering sum, difference, product, and inverse for all boundary
cases), which exceeds the theoretical bound for 4-limb operations.
\end{remark}

\section{Case Studies: Defects Invisible to Conventional Testing}
\label{sec:casestudies}

We present four defects from our ledger that demonstrate the
methodology catching bugs that the implementation's own test suite
could not detect.

\subsection{Case Study 1: Out-of-Bounds Memory Access (CA-R119)}

\subsubsection{The Defect}

The Vesta-domain Poseidon permutation (a separate implementation from
the Pallas twin) wrapped its full-round phases in a tripling loop:
\begin{lstlisting}[language=C++]
for (int r = 0; r < 4; ++r) {
    for (int k = 0; k < 3; ++k) {
        s.s0 = fe_add(s.s0, RC[k++]);
        s.s1 = fe_add(s.s1, RC[k++]);
        s.s2 = fe_add(s.s2, RC[k++]);
        // sbox, mix
    }
}
\end{lstlisting}
This consumed $4 \times 3 \times 3 = 36$ round constants per full phase,
for a total of $36 + 56 + 36 = 128$ reads against an 80-entry table.

\subsubsection{Why Self-Tests Missed It}

The implementation's own test suite verified: (1) domain separation,
(2) determinism, (3) sponge roundtrip, and (4) avalanche effect. All
four are \emph{self-consistency} tests: they verify that the output is
consistent with the input, but do not verify that the output is
\emph{correct}. The out-of-bounds reads produced deterministic (if
incorrect) values, so all self-consistency tests passed.

\subsubsection{How the Methodology Caught It}

The bilingual arithmetic cross-check computed the expected constraint
count: $\mathit{RF}/2 \times 3 + \mathit{RP} + \mathit{RF}/2 \times 3
= 12 + 56 + 12 = 80$ round constants. The loop structure consumed 128.
The discrepancy ($128 \neq 80$) was detected by pure arithmetic,
before any code was executed.

\subsubsection{The Fix}

The tripling loop was removed, reducing the consumption to
 $4 \times 3 + 56 + 4 \times 3 = 80$ round constants. The
differential probe confirmed that the output changed after the fix,
proving the old outputs were incorrect.

\subsection{Case Study 2: Wrong-Power Invariant Checker (CA-R121)}

\subsubsection{The Defect}

The on-curve check for Jacobian coordinates evaluated
 $Y^2 = X^3 + b \cdot Z^3$ instead of the correct
 $Y^2 = X^3 + b \cdot Z^6$. The Jacobian coordinate transformation
requires $Z^6$ (since $x = X/Z^2$ and $y = Y/Z^3$, substitution
gives $Y^2/Z^6 = X^3/Z^6 + b$, hence $Y^2 = X^3 + b \cdot Z^6$).

\subsubsection{Why Self-Tests Missed It}

The only point ever tested with the on-curve check was the generator,
which has $Z = 1$. For $Z = 1$, both $Z^3$ and $Z^6$ equal 1, so the
wrong power was invisible. The test suite did not include any point
with $Z \neq 1$.

\subsubsection{How the Methodology Caught It}

The multifold circuit's output points had $Z \neq 1$ (from the Jacobian
accumulation in scalar multiplication). When the on-curve check was
applied to these points, the wrong power produced a false negative.
The affine coordinates (computed by a correct to\_affine function)
showed the point WAS on the curve, while the on-curve check said it
was not — a contradiction that convicted the checker.

\subsubsection{The Fix}

The $Z^3$ was replaced with $Z^6$. A regression test was added to
test\_step26 that applies the on-curve check to every scalar-multiplication
output (guaranteeing $Z \neq 1$ shapes).

\subsection{Case Study 3: Bit-Order Convention Mismatch (CA-R133)}

\subsubsection{The Defect}

The C++ implementation of the $\mathit{eq}$ table (used in the
WHIR-class polynomial commitment) bound the first $\tau$ coordinate
to the \emph{high} bit of the table index, while the Python oracle
bound it to the \emph{low} bit. This is a convention choice with no
inherently correct answer — but the two implementations chose
differently.

\subsubsection{How the Methodology Caught It}

An orientation probe evaluated the eq-table with both orderings and
compared against the golden vector. The reversed ordering matched
the golden bit-for-bit; the as-coded ordering did not. The probe
produced a single bit of evidence (reversed == golden) that convicted
the C++ convention.

\subsubsection{The Fix}

The C++ eq\_table was changed to reverse the bit binding. The law
was recorded: ``Bit-order in a fresh implementation is a contract,
not a default — the golden pins it.''

\subsection{Case Study 4: Double Canonicalization (CA-R132)}

\subsubsection{The Defect}

The WHIR wrap's commitment computation pre-canonicalized values
before passing them to a function that canonicalizes internally:
\begin{lstlisting}[language=C++]
// WRONG: double canonicalization
fp::fe_to_le_bytes(fp::fe_to_canonical(P.C), cb);
\end{lstlisting}
The function \texttt{fe\_to\_le\_bytes} already calls
\texttt{fe\_to\_canonical} internally. The outer call applied
the Montgomery reduction a second time, producing
 $\mathsf{Canon}(\mathsf{Canon}(x)) = \mathsf{Canon}(x) \cdot R^{-1}$,
which is garbage.

\subsubsection{How the Methodology Caught It}

A three-way consistency check: the commitment computed via the direct
path (single canonicalization, correct) vs.\ the commitment computed
via the wrapper (double canonicalization, incorrect). The two
disagreed on the same input, proving that one path was wrong.

\subsubsection{The Fix}

The redundant outer canonicalization was removed from all four
affected call sites. A law was recorded: \emph{a helper's internal
normalization is part of its contract; the caller must not
pre-normalize}.

\section{The System}
\label{sec:system}

We apply the methodology to a complete cryptographic proof system:
a Layer-1 protocol architecture with threshold-encrypted mempool,
folding-based recursive state transitions, and succinct epoch proofs.

\subsection{The Module Inventory}

\begin{table}[h]
\centering
\caption{Implementation modules and their verification status}
\begin{tabular}{lll}
\toprule
\textbf{Module} & \textbf{Description} & \textbf{Constraints} \\
\midrule
fe.hpp & Pallas field $\mathbb{F}_p$ arithmetic & 516 golden cases \\
fev.hpp & Vesta field $\mathbb{F}_q$ arithmetic & 516 golden cases \\
poseidon.hpp & Pallas-domain Poseidon-3 hash & 80 RC, 9 MDS \\
poseidon\_v.hpp & Vesta-domain Poseidon-3 hash & 80 RC, 9 MDS \\
g2v.hpp & Vesta curve operations & 24 golden triples \\
g1p.hpp & Pallas curve operations & 24 golden triples \\
pedersen.hpp & Homomorphic Pedersen (both curves) & 8 bases \\
whir.hpp & WHIR-class succinct opening & 8 OOD points \\
mfold.hpp & HyperNova multifold core & 24 golden triples \\
cfold.hpp & Commitment folding & identity check \\
fexec\_circuit.hpp & $F_{\text{exec}}$ as R1CS & 8 rows \\
poseidon\_r1cs.hpp & Poseidon R1CS gadget & 1088 rows \\
epoch\_golden.hpp & The epoch chain + $\pi_E$ & 1600 bytes \\
\bottomrule
\end{tabular}
\end{table}

\subsection{The $\pi_E$ Receipt}

The system produces a recursive state transition proof $\pi_E$:
\begin{table}[h]
\centering
\caption{$\pi_E$ size receipt at golden scale}
\begin{tabular}{lr}
\toprule
\textbf{Component} & \textbf{Size (bytes)} \\
\midrule
T1 (opening transcript) & 544 \\
T2 (batched fold transcript) & 768 \\
OOD (8 points + commitment) & 288 \\
\midrule
\textbf{Total $\pi_E$} & \textbf{1600} \\
\textbf{Budget} & 73728 \\
\textbf{Utilization} & 2.2\% \\
\bottomrule
\end{tabular}
\end{table}

The proof is verified without the witness (verifier independence)
in 118\,ms on ARM64 (Termux, phone-class hardware).

\subsection{Production-Scale Projection}

At $k = 476$ decree entries with the full Poseidon permutation
($r_f/2 = 4$ full half-rounds, $r_p = 56$ partial rounds):

\begin{table}[h]
\centering
\caption{Production-scale constraint counts}
\begin{tabular}{lr}
\toprule
\textbf{Component} & \textbf{Constraints} \\
\midrule
 $f_{\text{exec}}$ transition (8 rows $\times$ 476) & 3{,}808 \\
Poseidon hash (1{,}088 $\times$ 3 $\times$ 476) & 1{,}553{,}664 \\
SMT opening (5{,}120 $\times$ 2 $\times$ 476) & 4{,}874{,}240 \\
\midrule
\textbf{Total} & \textbf{6{,}431{,}712} \\
\textbf{10$^6$ target} & \textbf{Exceeded ($6.4\times$)} \\
\bottomrule
\end{tabular}
\end{table}

\section{Results}
\label{sec:results}

\subsection{Defect Inventory}

The methodology caught 162 defects across the following classes:

\begin{table}[h]
\centering
\caption{Defect classes and counts}
\begin{tabular}{lrr}
\toprule
\textbf{Class} & \textbf{Count} & \textbf{Example} \\
\midrule
Arithmetic (carry, overflow, borrow) & 28 & CA-R48 \\
Convention (bit-order, endianness) & 22 & CA-R133 \\
Domain (Montgomery vs.\ canonical) & 18 & CA-R148 \\
Constraint algebra (R1CS row construction) & 15 & CA-R143 \\
API misuse (wrong function, wrong power) & 12 & CA-R121 \\
Structural (loop bounds, index spaces) & 10 & CA-R119 \\
Out-of-bounds memory access & 3 & CA-R119 \\
Emitter formatting (C++ syntax from generator) & 8 & CA-R127 \\
Infrastructure (gate, build system, git) & 6 & CA-R151 \\
Other & 40 & CA-R85 \\
\midrule
\textbf{Total} & \textbf{162} & \\
\bottomrule
\end{tabular}
\end{table}

\subsection{Defects Invisible to Conventional Testing}

Of the 162 defects, we identify 4 that were \emph{structurally
invisible} to the implementation's own test suite:

\begin{enumerate}
\item \textbf{CA-R119}: Out-of-bounds memory access — the test suite
verified self-consistency but not correctness (Section~\ref{sec:casestudies}).
\item \textbf{CA-R121}: Wrong-power on-curve check — the only tested
point had $Z=1$, masking the error (Section~\ref{sec:casestudies}).
\item \textbf{CA-R133}: Bit-order mismatch — a convention choice with
no self-evident correct answer.
\item \textbf{CA-R132}: Double canonicalization — the same helper was
used in both the correct and incorrect paths, making the error
self-consistent.
\end{enumerate}

These defects would have survived conventional unit testing, integration
testing, fuzzing, and code review. They were caught exclusively by the
bilingual cross-examination and the arithmetic consistency checks.

\subsection{The Gate's Own Integrity}

The verification gate itself was audited during development. Six
integrity defects were found and fixed (CA-R151 through CA-R157):
build failures that fell through to stale binaries, golden-filename
collisions, stale generator dependencies, and shell syntax errors
in the gate script itself. Each defect was converted to a law,
and the gate now verifies its own integrity as part of every run.

\section{Comparison with Prior Work}

\begin{table}[h]
\centering
\caption{Comparison with existing verification approaches}
\begin{tabular}{lccc}
\toprule
\textbf{Approach} & \textbf{Bilingual} & \textbf{Golden} & \textbf{Ledger} \\
\midrule
Unit tests & \checkmark & & \\
Property-based testing & & \checkmark & \\
Formal verification (EasyCrypt, F*) & & \checkmark & \checkmark \\
Reference-implementation diff & \checkmark & & \\
\textbf{This work} & \checkmark & \checkmark & \checkmark \\
\bottomrule
\end{tabular}
\end{table}

The key differentiator is the \emph{combination}: bilingual
cross-examination (two independent implementations) + golden vectors
(every intermediate value pinned) + the append-only ledger (every
defect owned). No prior work combines all three.

\section{Limitations}

\begin{enumerate}
\item \textbf{Shared-model risk}: if both implementations share a
design error (e.g., both use the same wrong bit-order convention),
the cross-examination will not detect it. This is mitigated by
deriving the convention from a reference implementation (e.g., the
pasta-curves crate) and by CA-R37's law: \emph{trust the goldens,
not the story}.
\item \textbf{Oracle correctness}: the Python oracle is itself
software and may contain bugs. However, it uses arbitrary-precision
integers (no fixed-width arithmetic), eliminating an entire class
of implementation bugs.
\item \textbf{Coverage}: the golden vectors test specific inputs;
they do not prove correctness for all inputs. The gap is bridged
by the structural coverage analysis (Theorem~\ref{thm:coverage})
and by the pre-emit self-check law (CA-R126).
\end{enumerate}

\section{Conclusion and Future Work}

We presented a methodology for verifying cryptographic implementations
through bilingual cross-examination, golden-vector pinning, and an
append-only decision ledger. Applied to a complete cryptographic proof
system, the methodology caught 162 defects — including four classes
invisible to conventional testing. The resulting system produces a
1600-byte recursive state transition proof verifiable without the
witness.

Future work includes: (1) formal verification of the core arithmetic
using EasyCrypt or F*, (2) extension to the full production circuit
(the complete SMT depth, the complete Poseidon round count), and
(3) external security audit by an independent firm.

\section{Artifacts}

All artifacts are available at:
\url{https://github.com/sarinsk629-blip/hsma-core}

\begin{itemize}
\item \texttt{include/hsma/}: 41 implementation headers
\item \texttt{conformance/}: 34 test binaries
\item \texttt{scripts/gate.sh}: the verification gate
\item \texttt{docs/DECISIONS.md}: the append-only decision ledger
(234 decisions, 163 laws)
\item \texttt{tools/}: 20+ diagnostic probes
\end{itemize}

The verification pipeline is reproducible: clone the repository,
run \texttt{./scripts/gate.sh}, and observe the GATE GREEN verdict.

\end{document}
